\documentclass[12pt]{article}
\usepackage[margin=1in]{geometry}
\usepackage{titlesec}
\usepackage{hyperref}
\usepackage{booktabs}
\usepackage{array}
\usepackage{parskip}

\title{Snapshot Objectives\\
\large QTC Workflow Queue (WFQ) Administrator Application}
\author{Course Team --- CS3338}
\date{Version 1.0 \\ Snapshot 1 Start}

\begin{document}

\maketitle
\newpage

\tableofcontents
\newpage

% ─────────────────────────────────────────────
\section{Snapshot 1 — Start Objective}
% ─────────────────────────────────────────────

\subsection{Overview}
Snapshot 1 marks the beginning of our course team's engagement with the QTC Workflow Queue
(WFQ) Administrator Application project. This project is being used as a template to
demonstrate proficiency in a set of software engineering tools: GitHub, Docker, Jira,
TestRail, and LaTeX.

The original SRS and SDD documents were produced by the QTC project team (Leidos QTC Health
Services). Our team's role is to build on top of this existing foundation --- adapting the
documentation into LaTeX, setting up the development toolchain, and adding new features and
testing artifacts across each snapshot.

Snapshot 1 covers the first two months of the six-month project timeline and focuses on
toolchain setup and establishing the base project structure.

\subsection{Goals}

\begin{itemize}
    \item \textbf{Repository Setup (GitHub):} Initialize the GitHub repository with the
    required folder structure (\texttt{/docs}, \texttt{/testrail}), a README user manual,
    and all required LaTeX documents.

    \item \textbf{Documentation (LaTeX):} Convert the existing SRS and SDD documents into
    \texttt{.tex} format. Create the Design Specification and this Snapshot Objectives
    document from scratch.

    \item \textbf{Containerization (Docker):} Create a \texttt{docker-compose.yml} file
    that defines the application's services (frontend/backend and database), based on the
    technology stack defined in the SDD.

    \item \textbf{Project Management (Jira):} Set up a Jira project with Sprint 1 tasks
    derived from the existing requirements. Add known risk/bug tickets that will be used
    in TestRail test cases.

    \item \textbf{Base Application Framework:} Initialize the ASP.NET Core MVC project,
    configure Entity Framework Core with the existing WFQ SQL Server database, and
    implement the three-layer architecture (Presentation, Business Logic, Data Access).

    \item \textbf{Authorization Module:} Implement the URL query string flag-based
    authorization mechanism (\texttt{?is\_admin=true}).

    \item \textbf{Workflow Dashboard:} Build the core Workflow Management Dashboard,
    including a paginated, searchable, and sortable list of all existing workflows.

    \item \textbf{Workflow CRUD:} Implement the ability for administrators to create and
    update workflows, including all required field validation.
\end{itemize}

\subsection{Task Delegation}

\textit{Replace the names below with your actual course team members and their assigned
responsibilities.}

\begin{center}
\begin{tabular}{>{\bfseries}p{4cm} p{9cm}}
\toprule
Team Member & Responsibility \\
\midrule
(Member 1) & Repository setup, README, GitHub structure \\
(Member 2) & SRS.tex conversion and formatting \\
(Member 3) & SDD.tex conversion and formatting \\
(Member 4) & DesignSpec.tex, UI component breakdown \\
(Member 5) & docker-compose.yml, Docker service definitions \\
(Member 6) & Jira project setup, Sprint 1 task creation \\
(Member 7) & Authorization Module implementation \\
(Member 8) & Workflow Dashboard UI \\
(Member 9) & Workflow Create/Edit functionality \\
(Member 10) & SnapshotObjectives.tex, Workflow Diagram \\
\bottomrule
\end{tabular}
\end{center}

\subsection{Base Dependencies}

\begin{itemize}
    \item \textbf{Framework:} ASP.NET Core MVC (.NET 7)
    \item \textbf{ORM:} Entity Framework Core 7
    \item \textbf{Database:} Microsoft SQL Server 2019+
    \item \textbf{Web Server:} Internet Information Services (IIS) 10.0
    \item \textbf{IDE:} Visual Studio 2022
    \item \textbf{Source Control:} Git, hosted on GitHub
    \item \textbf{NuGet Packages:}
    \begin{itemize}
        \item \texttt{Microsoft.EntityFrameworkCore.SqlServer}
        \item \texttt{Microsoft.EntityFrameworkCore.Tools}
        \item \texttt{Microsoft.AspNetCore.Authorization}
    \end{itemize}
\end{itemize}

% ─────────────────────────────────────────────
\section{Snapshot 2 — 1st Checkpoint Objective}
% ─────────────────────────────────────────────

\subsection{Reflection on Snapshot 1}
Snapshot 1 established the full foundation of the project. The repository was structured on
GitHub with all required LaTeX documents (SRS, SDD, Design Specification, and this Snapshot
Objectives document). The base ASP.NET Core MVC application was initialized with Entity
Framework Core connected to the WFQ SQL Server database. The Authorization Module
(\texttt{?is\_admin=true} query string flag) was implemented, and the core Workflow
Management Dashboard was built, including paginated workflow listing, search, sort, workflow
creation, and workflow update functionality. Jira Sprint 1 was configured with tasks and
known risk tickets. The Docker \texttt{docker-compose.yml} scaffold was also created.

All Snapshot 1 goals were completed as planned.

\subsection{New Feature: Email Notification System}
Snapshot 2 introduces the \textbf{Email Notification System} --- a new service that
automatically sends email alerts to relevant users when key workflow events occur. This
feature was identified in the original SRS (Section 2.9) as a planned future addition
(Version 2.0) and is now being implemented as part of this course project.

\subsection{Goals}

\begin{itemize}
    \item \textbf{Notification Service (Backend):} Implement a new
    \texttt{INotificationService} interface and \texttt{EmailNotificationService} class
    within the Business Logic Layer. This service will be responsible for composing and
    sending email notifications using SMTP.

    \item \textbf{Trigger Points:} Integrate notification triggers into the existing
    \texttt{WorkflowAdminService} for the following events:
    \begin{itemize}
        \item A task is assigned or reassigned to a user.
        \item A task is marked as completed.
        \item A new workflow step is created.
    \end{itemize}

    \item \textbf{Email Templates:} Create simple, readable plain-text email templates for
    each notification type. Templates will include the workflow name, step name, task
    details, and a direct link to the relevant workflow in the application.

    \item \textbf{Configuration:} Add SMTP configuration settings to
    \texttt{appsettings.json} (host, port, sender address). Credentials will be stored
    securely and not committed to source control.

    \item \textbf{Documentation Updates:} Update the SRS (new functional requirements),
    SDD (new service module section), and Design Specification (notification behavior per
    page) to reflect the new feature.

    \item \textbf{Jira Sprint 2:} Create a new sprint with tasks covering backend service
    implementation, template creation, configuration, and testing.

    \item \textbf{TestRail Snapshot 2:} Create and execute a test run in TestRail covering
    the notification feature. Export the run summary report and add it to the
    \texttt{/testrail} folder.
\end{itemize}

\subsection{New Dependencies}

\begin{itemize}
    \item \textbf{NuGet Package:} \texttt{MailKit} (pinned version) --- a cross-platform
    .NET library for sending email via SMTP. Chosen over \texttt{System.Net.Mail} for
    better async support and testability.
    \item \textbf{SMTP Server:} A configured SMTP relay (e.g., Gmail SMTP, SendGrid, or a
    local MailHog instance for development/testing).
\end{itemize}

\subsection{Task Delegation}

\textit{Replace the names below with your actual course team members and their assigned
responsibilities for Snapshot 2.}

\begin{center}
\begin{tabular}{>{\bfseries}p{4cm} p{9cm}}
\toprule
Team Member & Responsibility \\
\midrule
(Member 1) & INotificationService interface and EmailNotificationService implementation \\
(Member 2) & SMTP configuration, appsettings.json updates, secrets management \\
(Member 3) & Email template design (task assigned, task completed, step created) \\
(Member 4) & Integration of notification triggers into WorkflowAdminService \\
(Member 5) & SRS updates (new functional requirements section) \\
(Member 6) & SDD updates (new Notification Service module section) \\
(Member 7) & Design Spec updates (notification behavior per page) \\
(Member 8) & Jira Sprint 2 setup and task creation \\
(Member 9) & TestRail test case creation and test run execution \\
(Member 10) & Workflow Diagram update, SnapshotObjectives.tex Snapshot 2 section \\
\bottomrule
\end{tabular}
\end{center}

% ─────────────────────────────────────────────
\section{Snapshot 3 — 2nd Checkpoint Objective}
% ─────────────────────────────────────────────

\subsection{Reflection on Snapshot 2}
Snapshot 2 introduced the Email Notification System, a new feature that extends the WFQ
Administrator Application with automated email alerts for key workflow events (task assigned,
task completed, step created). The \texttt{INotificationService} interface and
\texttt{EmailNotificationService} implementation were added to the Business Logic Layer using
the \textbf{MailKit} NuGet library. SMTP configuration was added to
\texttt{appsettings.json} with credentials managed via environment variables. Notification
triggers were integrated into the existing \texttt{WorkflowAdminService}. The SRS, SDD, and
Design Specification were updated to reflect the new feature. Jira Sprint 2 was configured
and TestRail Snapshot 2 test cases were created and executed.

All Snapshot 2 goals were completed as planned.

\subsection{New Feature: Workflow Export / Report Generation}
Snapshot 3 introduces the \textbf{Workflow Export and Report Generation} feature --- a new
capability that allows administrators to export a workflow's full details (steps, tasks, and
history) as a downloadable \textbf{PDF} or \textbf{CSV} file directly from the Workflow
Details page.

This feature was not present in the original project and is a new addition by the CS3338
course team. It introduces a new \texttt{IExportService} component to the Business Logic
Layer and adds export action endpoints to the \texttt{WorkflowsController}.

\subsection{Goals}

\begin{itemize}
    \item \textbf{Export Service (Backend):} Implement a new \texttt{IExportService}
    interface and \texttt{WorkflowExportService} class within the Business Logic Layer.
    This service will be responsible for generating PDF and CSV representations of a
    workflow's data.

    \item \textbf{PDF Export:} Generate a formatted PDF report containing the workflow
    name, description, status, a table of all steps (StepNumber, StepName, Description),
    and a table of all tasks (AssignedTo, Status, Step). Use the \textbf{QuestPDF} NuGet
    library for PDF generation.

    \item \textbf{CSV Export:} Generate a flat CSV file containing all tasks for a selected
    workflow, with columns: WorkflowName, StepName, StepNumber, TaskID, AssignedTo, Status,
    StartDate, CompletionDate.

    \item \textbf{UI Integration:} Add ``Export PDF'' and ``Export CSV'' buttons to the
    Workflow Details page. Clicking either button triggers a file download in the browser
    without navigating away from the page.

    \item \textbf{Documentation Updates:} Update the SRS (new functional requirements),
    SDD (new Export Service module section), and Design Specification (new export buttons
    on the Workflow Details page) to reflect the new feature.

    \item \textbf{Jira Sprint 3:} Create a new sprint with tasks covering backend service
    implementation, PDF/CSV generation, UI button integration, and testing.

    \item \textbf{TestRail Snapshot 3:} Create and execute a test run in TestRail covering
    the export feature. Export the run summary report and add it to the
    \texttt{/testrail} folder.
\end{itemize}

\subsection{New Dependencies}

\begin{itemize}
    \item \textbf{NuGet Package:} \texttt{QuestPDF} (pinned version) --- a modern,
    fluent .NET library for generating PDF documents. Chosen for its clean API, active
    maintenance, and suitability for server-side rendering without external dependencies.
    \item \textbf{No additional packages required for CSV} --- CSV generation will be
    handled using the built-in \texttt{System.Text} and \texttt{StringBuilder} classes to
    avoid unnecessary dependencies.
\end{itemize}

\subsection{Task Delegation}

\textit{Replace the names below with your actual course team members and their assigned
responsibilities for Snapshot 3.}

\begin{center}
\begin{tabular}{>{\bfseries}p{4cm} p{9cm}}
\toprule
Team Member & Responsibility \\
\midrule
(Member 1) & IExportService interface and WorkflowExportService implementation \\
(Member 2) & PDF report layout and generation using QuestPDF \\
(Member 3) & CSV generation logic using StringBuilder \\
(Member 4) & Export controller actions and file download response handling \\
(Member 5) & UI integration --- Export PDF and Export CSV buttons on Workflow Details page \\
(Member 6) & SRS updates (new functional requirements section) \\
(Member 7) & SDD updates (new Export Service module section) \\
(Member 8) & Design Spec updates (export buttons and behavior) \\
(Member 9) & Jira Sprint 3 setup and task creation \\
(Member 10) & TestRail test case creation, test run execution, and report export \\
\bottomrule
\end{tabular}
\end{center}

% ─────────────────────────────────────────────
\section{Snapshot 4 — Due Date Checkpoint}
% ─────────────────────────────────────────────

\textit{To be completed at Snapshot 4.}

\end{document}
